% =============================================================================
% ABSTRACT - v3 EMERGENT SOCIAL STRUCTURES FRAMING
% =============================================================================
% STATUS: Needs rewrite around cooperation literature + semantic priors
%
% FRAMING EVOLUTION:
%   v1: "Neorealism vs Constructivism" → Circular
%   v2: "Security Dilemma Escape" → Too narrow
%   v3: "Emergent Social Structures" → Captures full game richness
%
% NEW STRUCTURE (~100-150 words):
%   1. BACKGROUND: Cooperation among self-optimizers is well-studied (Axelrod)
%   2. LIMITATION: This cooperation is mechanistic - agents don't "understand" it
%   3. QUESTION: Do semantic priors enable qualitatively DIFFERENT structures?
%   4. APPROACH: Compare RL (trial-and-error) vs LLM (semantic reasoning),
%      SAME optimization goal, different reasoning architecture
%   5. MEASURES: Coalition formation, cooperation patterns, hierarchy, conflict
%   6. CONTRIBUTION: If structures differ → semantic priors matter for coordination
%
% TYPOS TO FIX:
%   - "comparitive" → "comparative"
%   - "Reinforcemnt" → "Reinforcement"
%
% EXAMPLE DIRECTION (don't copy, write in your voice):
%
%   "Cooperation among self-interested agents has been extensively studied.
%   Axelrod (1984) showed that optimizing agents can achieve cooperation
%   through repeated interaction. However, this cooperation emerges
%   mechanistically - agents discover stable equilibria without understanding
%   cooperation.
%
%   This thesis asks whether semantic reasoning enables qualitatively different
%   emergent social structures. We compare self-optimizing RL agents (trial-
%   and-error learning) with self-optimizing LLM agents (semantic reasoning)
%   in a multi-agent coordination game. Both pursue the same goal; they differ
%   in HOW they reason.
%
%   We measure coalition formation, cooperation patterns, hierarchy emergence,
%   and conflict dynamics. If LLM agents produce richer structures - more
%   stable coalitions, faster trust-building, emergent norms - this suggests
%   semantic priors enable coordination that pure optimization cannot access."
%
% KEY CHANGES FROM v2:
%   - Lead with cooperation literature (Axelrod), not security dilemma
%   - Frame as "what do semantic priors ADD?" not "can they escape dilemmas?"
%   - Emphasize MULTIPLE emergent structures, not just cooperation/conflict
%   - Position as extending classic cooperation studies into semantic reasoning era
%
% OPTIONAL: Brief IR mention as application domain, but NOT the main framing
% =============================================================================

\begin{abstract}
Thirty years since A. Wendt's Constructivism challenged K. Waltz's Neorealism both International Relations (IR) theories are acknowledged as major pillars to the field. Although Agent-based modelling (ABM) has long made its way into IR, direct comparisons in the same simulation between neorealism and constructivism have been absent. Recent advancements in Artificial Intelligence allow for a comparitive study as constructivist actor macro-behaviour can be modelled through Large-Language Model (LLM) agents. Reinforcemnt Learning (RL) agents represent materialistic Neorealist actors.

My thesis will explore under what material conditions the representations of Neorealistic actors and Constructivist actors will produce divergent emergent social structures.

\end{abstract}
